\section{Dogfood Strategy}

Atlas dogfood is not a demo. It is the path by which Kungfu learns what real
agent work management requires.

Atlas already contains mission/go semantics, worktree isolation, task ledgers,
goal registries, review records, dogfood queues, multi-agent handoff, and
external repository coordinates. Many of these mechanisms are hand-built from
Git and JSONL. They work, but they also expose the cost of managing runtime
facts without a product designed for that purpose.

Kungfu should absorb the lessons, not merely import the terminology. The
sequence is:

\begin{enumerate}[leftmargin=*]
  \item import Atlas mission/go facts into Kungfu so the GUI can display them;
  \item let agents write runtime facts, traces, friction, and task state into
  Kungfu while Atlas remains a backup source during the transition;
  \item identify which Atlas semantics are personal workflow choices and which
  belong in the generic Kungfu core;
  \item rebuild future Atlas mechanisms around the generic core where that
  reduces drift and token-heavy coordination;
  \item keep dogfood pressure high enough that Kungfu must become useful in
  daily work before public launch.
\end{enumerate}

The dogfood standard is load-bearing: if Kungfu cannot manage the work that is
building Kungfu, the product thesis has not been proven.
